\documentclass[11pt,a4paper]{article}

\usepackage[utf8]{inputenc}
\usepackage[T1]{fontenc}
\usepackage{lmodern}
\usepackage[margin=2.5cm]{geometry}
\usepackage{microtype}
\usepackage{amsmath,amssymb,mathtools}
\usepackage{booktabs,array,tabularx,longtable,threeparttable}
\usepackage{enumitem}
\usepackage{xcolor}
\usepackage{listings}
\usepackage{caption}
\usepackage{float}
\usepackage[hidelinks]{hyperref}
\usepackage{xurl}
\emergencystretch=2.5em

\setlist{itemsep=2pt,topsep=4pt}
\captionsetup{font=small,labelfont=bf}
\setlength{\parskip}{3pt}

\newcommand{\code}[1]{\texttt{\detokenize{#1}}}
\DeclareRobustCommand{\src}[1]{{\small\nolinkurl{#1}}}
\newcommand{\lbl}[1]{\textsc{#1}}

\definecolor{codebg}{gray}{0.96}
\lstset{
  basicstyle=\ttfamily\footnotesize,
  backgroundcolor=\color{codebg},
  frame=single, rulecolor=\color{gray!50},
  columns=fullflexible, keepspaces=true,
  breaklines=true, breakatwhitespace=false,
  showstringspaces=false, aboveskip=6pt, belowskip=6pt,
  xleftmargin=2pt, xrightmargin=2pt
}

\title{How ConflictScore, RAGChecker and Q-CARE Are Implemented and Scored\\[4pt]
\large A methodology and implementation guide for the conflicting-evidence evaluator benchmark}
\author{Factuality and Faithfulness Evaluation project}
\date{19 November 2026 (implementation 1.0.0)}

\begin{document}
\maketitle

\begin{abstract}
\noindent
This guide explains the three published evaluators that the project runs: ConflictScore, RAGChecker and Q-CARE. For each one it covers the original paper and benchmark, the exact scoring formulas, and the path a single example takes through our Python code. It ends with a step-by-step recomputation of real C1 results. Every statement about our implementation was checked against the source code in \src{src/benchmark/}, the pinned upstream code in \src{external/} (or the installed \code{ragchecker}/\code{refchecker} packages), and the stored outputs in \src{outputs/runs/} and \src{outputs/results/}. Statements about the papers come from the papers themselves, or from the official repository README where the paper text was not consulted. The last section lists every discrepancy found between the published methodology and the version run here. Results are summarised in \src{outputs/reports/evaluation_results_brief.pdf}; a two-page overview is \src{outputs/reports/implementation_summary.pdf}.
\end{abstract}

\tableofcontents
\newpage

%=====================================================================
\section{The shared experimental setting}
\label{sec:setting}

\subsection{The controlled C1 example}

All worked examples use scenario \lbl{C1}, the reference condition. The scenario is built by \code{build_scenarios()} in \src{src/benchmark/data_factory.py:161}. It has two retrieved documents about a fictional person. Both have the same assigned reliability, both are stated without hedging, and they disagree only on the day. \code{render_doc(..., "attributed")} (\src{data_factory.py:218}) renders them as the evaluators see them:

\begin{lstlisting}
ctx0: [Source: Archive A] Alex Meron was born on 3 March 1990.
ctx1: [Source: Registry B] Alex Meron was born on 17 March 1990.
\end{lstlisting}

The assigned reliability values (\code{high}=0.9) stay in the scenario metadata and never reach any evaluator. All three implementations accept only text (\code{TEXT_ONLY_METADATA}, \src{evaluators/base.py:71}). The final C1 run uses two question wordings:
\lbl{Q\_YEAR} ``Which year was Alex Meron born?'' and
\lbl{Q\_DATE} ``What exact date was Alex Meron born?''.
The candidate answers are fixed strings (\code{RESPONSES}, \src{data_factory.py:27}):

\begin{table}[h]
\centering\small
\begin{tabularx}{\textwidth}{llX}
\toprule
ID & Label & Text\\
\midrule
R1 & EXACT-A & Alex Meron was born on 3 March 1990.\\
R2 & MONTH & Alex Meron was born in March 1990.\\
R3 & YEAR & Alex Meron was born in 1990.\\
R4 & VAGUE & The exact date of birth is uncertain.\\
R5 & ATTRIBUTED-CONFLICT & Archive A gives 3 March 1990, while Registry B gives 17 March 1990. Both agree on March 1990, but the precise day is disputed.\\
\bottomrule
\end{tabularx}
\caption{The five main candidate answers. They are generated by no model; they are fixed test inputs.}
\label{tab:responses}
\end{table}

\subsection{End-to-end control flow}

A run is started by \code{python -m benchmark.cli run}. In order, \code{cmd_run} (\src{src/benchmark/cli.py:45}):
\begin{enumerate}
\item builds the judge from a spec string such as \code{codex:gpt-6-luna:low} (\code{make_judge}, \src{cli.py:17});
\item expands the scenario $\times$ question grid with \code{grid()} (\src{data_factory.py:195}), which also attaches the evidence-grounded reference used by RAGChecker (\code{golden_reference()}, \src{data_factory.py:116});
\item creates one \emph{item} for each (instance, response), with \code{item_id} = \code{C1|Q_DATE|R5}, \code{query}, \code{answer}, \code{contexts} (the rendered strings above) and \code{reference};
\item passes the items to each adapter's \code{evaluate_many()}. RAGChecker is called once per scenario$\times$question, so the five answers to one question share one official \code{evaluate()} call;
\item appends one JSON row per (evaluator, item) to \src{<run>/results.jsonl} via \code{emit()} (\src{cli.py:72}). Each row holds the native metric values, the full raw intermediate output (\code{raw}), the adaptation list and the judge description.
\end{enumerate}
Every adapter returns the same \code{EvaluationResult}/\code{MetricValue} structure (\src{evaluators/base.py:31-62}). No metric is rescaled or merged across evaluators.

\subsection{Where results are saved}

\begin{table}[H]
\centering\small
\begin{tabularx}{\textwidth}{lX}
\toprule
File & Contents\\
\midrule
\src{outputs/runs/main/results.jsonl} & Main run: 240 results (3 evaluators $\times$ 8 scenarios $\times$ 2 questions $\times$ 5 answers), all \lbl{OK}. Other runs in \src{outputs/runs/}: \src{ref_held_C1}, \src{reliability_in_text}, \src{judge_claude-sonnet}, \src{judge_claude-opus}.\\
\src{outputs/runs/*/judge_calls.jsonl} & Per-run judgment log (\code{call}, \code{cache_hit}, \code{error}) with the cache key; \src{run_meta.json} holds the command and judge configuration.\\
\src{outputs/results/judge_records.jsonl} & Exported full judge records: prompt, system text, response, token usage, timestamp, CLI version.\\
\src{cache/judge/<sha256>.json} & The content-addressed judge cache; every reported number can be replayed from it.\\
\src{outputs/results/comparison.csv} & Long table across all runs, produced by \code{report_data.comparison()} (\src{scripts/build_results.py}).\\
\src{outputs/results/claim_level_traces.json} & Condensed claims and labels for each result (\code{report_data.traces()}).\\
\src{outputs/results/judge_comparison_*.csv} & Cross-judge comparison of the same C1 inputs under Claude judges.\\
\src{outputs/results/judgment_provenance.csv} & For each run: number of cached judgments used, their generation dates and CLI versions, and new calls (always 0).\\
\bottomrule
\end{tabularx}
\caption{Output locations relevant to this guide.}
\end{table}

\subsection{The judge layer}
\label{sec:judge}

A \emph{judge} is the LLM that a published evaluator calls internally. The base class (\src{src/benchmark/judges/base.py}) exposes \code{chat(messages)}, \code{complete(prompt, system)} and \code{batch(prompts)}. \code{chat()} joins every \code{system} message into one system string and sends the single \code{user} message as the prompt (\src{base.py:97}). The production backend is \code{CodexJudge} (\src{judges/codex.py:87}), which runs \code{codex exec} under the user's ChatGPT login. The cross-judge backend \code{ClaudeJudge} (\src{judges/claude.py}) runs \code{claude -p}. For RAGChecker and Q-CARE each adapter replaces only the transport object that the official code calls (Table~\ref{tab:transport}); their prompts, parsers and scoring arithmetic are the upstream ones. ConflictScore's CS-C/CS-R have no official code, so the adapter sends the paper's prompts directly.

\begin{table}[h]
\centering\small
\begin{tabularx}{\textwidth}{p{2.4cm}>{\raggedright\arraybackslash}X>{\raggedright\arraybackslash}X}
\toprule
Evaluator & Official call site & Our substitute\\
\midrule
ConflictScore & none (no official CS-C/CS-R code) & paper prompts sent with \code{judge.complete} (\src{conflictscore.py:242})\\
RAGChecker & \code{custom_llm_api_func(prompts)}, a documented constructor hook & \code{RAGCheckerAdapter._llm} (\src{ragchecker.py:56}) $\rightarrow$ \code{judge.batch}\\
Q-CARE & \code{Backbone.generate(messages)} (local HF model) & \code{JudgeBackbone.generate} (\src{qcare.py:50}) $\rightarrow$ \code{judge.chat}\\
\bottomrule
\end{tabularx}
\caption{How each evaluator reaches the judge.}
\label{tab:transport}
\end{table}

%=====================================================================
\section{ConflictScore}
\label{sec:cs}

\subsection{A. Original paper and benchmark}

\begin{description}[leftmargin=0pt]
\item[Paper.] Siyi Liu, Aaron Halfaker, Dan Roth, Patrick Xia. \emph{ConflictScore: Identifying and Measuring How Language Models Handle Conflicting Evidence.} Proceedings of EMNLP 2026 (Main Conference), Budapest, pp.~58825--58845. \url{https://aclanthology.org/2026.emnlp-main.2589/}
\item[Code.] \url{https://github.com/microsoft/ConflictScore} (MIT), pinned at commit \code{a8868e5} (13 March 2026).
\item[Contribution type.] There are three contributions: a \emph{metric} (ConflictScore, with two aggregate measures CS-C and CS-R), a \emph{benchmark} (ConflictBench, built for the conflict-detection task), and two \emph{downstream studies}: benchmarking frontier LLMs, and using the score as feedback to improve TruthfulQA accuracy.
\item[ConflictBench.] ConflictBench contains 2{,}293 items (1{,}280 conflicting, 1{,}013 non-conflicting) taken from ContraQA (798), MacNoise-NQ (199), MacNoise-TQA (211), AmbigDocs (651) and ConflictingQA (434) (paper Table~1). The data are not released; the repository rebuilds them from the source datasets in four scripted steps.
\item[Judge models.] Main conflict-detection results (Table~2) use Gemini-3.5-Flash. Appendix~B repeats them with GPT-5.4-mini and Claude-haiku-4.5. Claims for ConflictBench are produced with GPT-4.1 (Appendix~A.1). The cost profile in the Limitations section uses Claude-haiku-4.5 via OpenRouter. \emph{Repository default}: \code{OPENAI_MODEL} = \code{gpt-4.1}, \code{temperature=0}, \code{max_completion_tokens=4096}.
\end{description}

\paragraph{Two different protocols in one paper.}
The paper describes two separate uses of the same labelling machinery. They should not be confused.
\begin{enumerate}
\item \emph{Response scoring (paper Section~3).} The response is decomposed into atomic claims (prompt in paper Table~10). Each claim is labelled against each document as supports, contradicts or irrelevant (Table~11), and the labels are aggregated into CS-C and CS-R. Sections~5--6 use this protocol.
\item \emph{Conflict detection on ConflictBench (Section~4).} Each item has \emph{one} claim, produced from the dataset's gold answer. The task is binary: does the claim have both supporting and contradicting documents? The repository's \src{src/} pipeline implements this protocol.
\end{enumerate}

\paragraph{How our experiment differs.}
The paper scores long model-generated reports (100 ConflictingQA items, three paragraphs each) or detects conflicts for one gold-derived claim per item. Our experiment scores five \emph{fixed, hand-written} answers against two documents that differ by one token (the day). No candidate answer is generated, there is no gold conflict label to calibrate against, and the paper's Section~4 classification metrics (precision, recall, F1 against gold labels) are not computed.

\subsection{B. Metric definitions and scoring logic}

\paragraph{Paper definitions.}
For a response with atomic claims $C=\{c_1,\dots,c_n\}$ and grounding documents $D$, let $D^+_i\subseteq D$ be the documents labelled as supporting $c_i$ and $D^-_i$ those labelled as contradicting it. The paper defines
\begin{align}
\text{conflicted}(c_i) &= \bigl[\,|D^+_i|>0 \;\wedge\; |D^-_i|>0\,\bigr],\\
\text{CS-C} &= \frac{1}{n}\sum_{i=1}^{n} \text{conflicted}(c_i), \\
\text{CS-R} &= \frac{1}{n}\sum_{i=1}^{n} \frac{|D^-_i|}{|D^+_i|+|D^-_i|}.
\end{align}
For both measures, \emph{lower is better}. The paper does not say how CS-R treats a claim with $|D^+_i|+|D^-_i|=0$.

\paragraph{What our code computes.}
The official repository does not contain CS-C or CS-R (verified against upstream HEAD). We implement them from the paper as mode \code{published_cs} (\src{evaluators/conflictscore_published.py}), using the prompts of Appendix~E word for word (Table~10: 45 words; Table~11: 342 words; only PDF typography such as line-break hyphenation and quotation marks is normalised). This is the project's only ConflictScore evaluation. The repository's \src{src/} pipeline (the ConflictBench detection protocol, which rewrites a gold answer as a single claim) and its example script use different prompts and scores and are not used (Deviations~\ref{d:cs-protocol} and~\ref{d:cs-prompts}).

\begin{table}[h]
\centering\small
\begin{tabularx}{\textwidth}{>{\raggedright\arraybackslash}p{6.4cm}Xc}
\toprule
Stored quantity & Definition & Better\\
\midrule
\code{cs_c} & CS-C over the Table~10 claims and Table~11 labels & lower\\
\code{cs_r} & CS-R; claims with $|D^+|+|D^-|=0$ excluded (undefined in the paper) & lower\\
\code{n_claims} & number of claims from the Table~10 decomposition & diagnostic\\
\code{n_conflicted_claims} & claims with $|D^+|>0$ and $|D^-|>0$ & diagnostic\\
\code{n_claims_without_relevant_evidence} & claims with no supporting or contradicting document (excluded from CS-R) & diagnostic\\
\bottomrule
\end{tabularx}
\caption{ConflictScore quantities stored for each answer. None of them uses a golden reference.}
\end{table}

\paragraph{LLM versus deterministic parts.}
\emph{LLM}: the claim decomposition and every claim--document label.
\emph{Deterministic}: parsing the claim list and the labels, counting, the per-claim ratio, the conflict rule and the averages over claims.

\subsection{C. What our code actually does}

\paragraph{Implementation} (\code{ConflictScoreAdapter._published}, \src{conflictscore.py:242}; prompts, parsers and aggregation in \src{conflictscore_published.py}).
For each answer:
\begin{enumerate}
\item \textbf{Decomposition (LLM).} One call with the Table~10 prompt (``Break down the following report into individual claims \ldots\ List each claim on a new line, starting with "Claims: "''), with the answer as the report. The prompt contains \emph{no question}.
\item \textbf{Claim list (deterministic).} The paper publishes no parser. Judges follow the instruction in two ways: ``Claims:'' before every claim (Codex, Sonnet) or one ``Claims:'' header followed by one claim per line (Opus). Both are accepted. If nothing parses, the whole answer is used as one claim and a warning is stored; this never happened in the reported runs.
\item \textbf{Evidence labels (LLM).} One call per (claim, document) with the Table~11 prompt, which also contains no question and asks for a JSON object with \code{"answer": "<SUPPORTS|CONTRADICTS|IRRELEVANT>"}.
\item \textbf{Label parsing (deterministic).} A line-for-line port of the official parser (\src{eval_all_calibration.py:104-140}): JSON \code{answer} field, regex fallback, \code{<answer>} tag fallback, default IRRELEVANT. Every label in the reported runs parsed to one of the three labels.
\item \textbf{Aggregation (deterministic).} CS-C and CS-R as defined in Section~B, with the convention for claims without relevant evidence.
\end{enumerate}
A judge error makes the answer \lbl{ERROR}; it is never counted as IRRELEVANT (the repository's API wrapper would do that). Because neither prompt contains the question, Q\_YEAR and Q\_DATE give identical results.

\subsection{D. Worked examples (C1, judge \code{gpt-6-luna})}

\begin{table}[h]
\centering\small
\begin{tabular}{l>{\raggedright\arraybackslash}p{7.2cm}ccc}
\toprule
Answer & Claims (Table~10) and labels (documents A, B) & CS-C & CS-R & excluded\\
\midrule
R1 & ``Alex Meron was born on 3 March 1990.'' (SUP, CON) & $1/1=1.00$ & $1/2=0.50$ & 0\\
R2, R3 & one claim each (SUP, SUP) & 0 & 0 & 0\\
R4 & ``The exact date of birth is uncertain.'' (CON, CON) & 0 & $2/2=1.00$ & 0\\
R5 & ``Archive A gives 3 March 1990.'' (SUP, IRR); ``Registry B gives 17 March 1990.'' (IRR, SUP); ``Both agree on March 1990.'' (IRR, SUP); ``The precise day is disputed.'' (IRR, IRR) & 0 & $\tfrac13(0+0+0)=0$ & 1\\
\bottomrule
\end{tabular}
\caption{ConflictScore claims and labels for C1 (identical for both questions; \src{outputs/runs/main/results.jsonl}).}
\label{tab:cs-trace}
\end{table}

\paragraph{Why R1 receives CS-C 1.00 and CS-R 0.50 on both questions.}
The single claim states the day 3. Archive A supports it and Registry B contradicts it, so $|D^+|=|D^-|=1$: the claim is conflicted (CS-C $=1/1$) and its contradiction ratio is $1/(1+1)=0.50$. The decomposition never sees the question, so the contested day is scored even when only the year was asked.

\paragraph{Why R4 is not conflicted but has CS-R 1.00.}
Each document gives an exact date, which the judge reads as contradicting ``the date is uncertain''. With $|D^+|=0$ and $|D^-|=2$ the claim is not conflicted under the paper's rule (CS-C $=0$), but its contradiction ratio is maximal ($2/2$).

\paragraph{Why R5 receives 0 and 0.}
Under the Table~11 prompt, the source that an attributed claim does \emph{not} name is IRRELEVANT, not contradicting, so no claim is conflicted and every defined ratio is 0. ``The precise day is disputed'' is IRRELEVANT for both documents; its ratio $0/0$ is undefined and the claim is excluded from CS-R. R5 therefore ties with R2 and R3. CS-C and CS-R only penalise; they cannot reward the acknowledgement of the conflict.

\paragraph{Same scores under the Claude judges.}
Claude Sonnet~5.5 and Opus~5.5 give exactly the same CS-C and CS-R for all five C1 answers, although their claim wording differs and Opus splits R2 into ``born in March'' and ``born in 1990'' (\src{outputs/results/judge_comparison_conflictscore_labels.csv}).

\paragraph{Effect of the CS-R convention.}
Every R5 has one claim without relevant evidence, and in C4 two. Counting such claims as 0 instead of excluding them would change two reported values, C4 R5 from 0.25 to 0.125 and NC R5 from 0.33 to 0.25, and no ranking. Results for C1--C5 and the controls are in \src{outputs/reports/evaluation_results_brief.pdf} and \src{outputs/results/conflictscore_published.csv}.

%=====================================================================
\section{RAGChecker}
\label{sec:rc}

\subsection{A. Original paper and benchmark}

\begin{description}[leftmargin=0pt]
\item[Paper.] Dongyu Ru, Lin Qiu, Xiangkun Hu, Tianhang Zhang, Peng Shi, Shuaichen Chang, Cheng Jiayang, Cunxiang Wang, Shichao Sun, Huanyu Li, Zizhao Zhang, Binjie Wang, Jiarong Jiang, Tong He, Zhiguo Wang, Pengfei Liu, Yue Zhang, Zheng Zhang. \emph{RAGChecker: A Fine-grained Framework for Diagnosing Retrieval-Augmented Generation.} NeurIPS 2024, Datasets and Benchmarks Track. arXiv:\href{https://arxiv.org/abs/2408.08067}{2408.08067}.
\item[Code.] \url{https://github.com/amazon-science/RAGChecker} (Apache-2.0). We run the installed packages \code{ragchecker==0.1.9} and \code{refchecker==0.2.17}. The installed \code{ragchecker} source is identical to \src{external/RAGChecker/ragchecker} (checked with \code{diff -rq}).
\item[Contribution type.] There are three contributions: an \emph{evaluation framework} with 11 diagnostic metrics; a \emph{benchmark} of 4{,}162 queries across 10 domains (ClapNQ 300; NovelQA 280; RobustQA Writing, Finance, Lifestyle, Recreation, Science and Technology 500 each; RobustQA BioASQ 511; KIWI 71); and a \emph{human meta-evaluation} on 280 pairwise instances.
\item[Original experiments.] Eight RAG systems: two retrievers (BM25, E5-Mistral) combined with four generators (GPT-4-turbo, Mixtral-8x7B, Llama3-8B, Llama3-70B). Each query has a human or dataset ground-truth answer.
\item[Judge models.] \emph{Paper}: Llama3-70B-Instruct as both claim extractor and checker. \emph{Code default}: \code{bedrock/meta.llama3-70b-instruct-v1:0} (\code{RAGChecker.__init__}). \emph{README examples}: \code{bedrock/meta.llama3-1-70b-instruct-v1:0} (Llama~3.1). The RefChecker class defaults (Claude-3-Sonnet on Bedrock) are overridden by RAGChecker.
\end{description}

\paragraph{How our experiment differs.}
RAGChecker was designed for queries with a single ground-truth answer. Our C1 question has no verified answer: the entity is fictional and the documents disagree. The project therefore supplies an \emph{evidence-grounded reference}, written before scoring (\code{GOLDEN_REFERENCES}, \src{data_factory.py:81}). It is not verified ground truth. For Q\_DATE it is
\begin{quote}\small
``The retrieved records disagree about Alex Meron's exact birth date. Archive A reports 3 March 1990, whereas Registry B reports 17 March 1990. Both agree on March 1990. The available evidence does not establish which day is correct.''
\end{quote}
For Q\_YEAR it is ``Alex Meron was born in 1990.'' Ten of the eleven metrics depend on this text.

\subsection{B. Metric definitions and scoring logic}

\paragraph{The three LLM-produced inputs.}
RefChecker's \code{LLMExtractor} turns a text into knowledge triplets (\emph{subject, predicate, object}). RAGChecker extracts response claims $m=\{m_1..m_p\}$ from the answer and ground-truth claims $g=\{g_1..g_q\}$ from the reference, and stores four entailment checks. The naming is \emph{reference}2\emph{claims}:
\begin{center}\small
\begin{tabular}{lll}
\toprule
Name & Claims checked & Checked against\\
\midrule
\code{answer2response}[$i$] & response claim $m_i$ & the reference text (one label)\\
\code{response2answer}[$j$] & reference claim $g_j$ & the response text (one label)\\
\code{retrieved2response}[$i$][$k$] & response claim $m_i$ & chunk $k$ (one label per chunk)\\
\code{retrieved2answer}[$j$][$k$] & reference claim $g_j$ & chunk $k$\\
\bottomrule
\end{tabular}
\end{center}
Each label is Entailment, Neutral or Contradiction. \code{to_bool} maps \emph{Entailment} to 1 and \emph{both} Neutral and Contradiction to 0 (\src{ragchecker/computation.py:7}). After that point RAGChecker no longer distinguishes ``contradicted'' from ``not mentioned''.

Write the Entailment indicators as
\begin{align*}
A_i&=\mathbb 1[\text{answer2response}_i=\text{Ent}], & B_j&=\mathbb 1[\text{response2answer}_j=\text{Ent}],\\
E_{ik}&=\mathbb 1[\text{retrieved2response}_{ik}=\text{Ent}], & G_{jk}&=\mathbb 1[\text{retrieved2answer}_{jk}=\text{Ent}],
\end{align*}
and let $K$ be the number of chunks. Chunk $k$ is \emph{relevant} if $\max_j G_{jk}=1$. The definitions below are transcribed from \src{computation.py} and agree with the paper's Appendix~B.

{\small\begin{longtable}{>{\raggedright\arraybackslash}p{4.4cm}p{6.2cm}p{1.5cm}p{2.2cm}}
\toprule
Metric & Formula (per item) & Better & Needs\\
\midrule
\endhead
\code{precision} & $\frac1p\sum_i A_i$ & higher & reference\\
\code{recall} & $\frac1q\sum_j B_j$ & higher & reference\\
\code{f1} & $\frac{2PR}{P+R}$ if $P,R>0$, else 0 & higher & reference\\
\code{claim_recall} & $\frac1q\sum_j \max_k G_{jk}$ & higher & ref.\ + chunks\\
\code{context_precision} & $\frac1K\sum_k \max_j G_{jk}$ & higher & ref.\ + chunks\\
\code{context_utilization} & $\dfrac{\sum_j \max_k G_{jk}\, B_j}{\sum_j \max_k G_{jk}}$ (0 if the denominator is 0) & higher & ref.\ + chunks\\
\code{noise_sensitivity_} \code{in_relevant} & $\frac1p\sum_i (1-A_i)\,\max_{k\,\text{rel}} E_{ik}$ & lower & ref.\ + chunks\\
\code{noise_sensitivity_} \code{in_irrelevant} & $\frac1p\sum_i (1-A_i)\max\limits_{k\,\text{irr}} E_{ik}(1-\max\limits_{k\,\text{rel}}E_{ik})$ & lower & ref.\ + chunks\\
\code{hallucination} & $\frac1p\sum_i (1-A_i)\,(1-\max_k E_{ik})$ & lower & ref.\ + chunks\\
\code{self_knowledge} & $\frac1p\sum_i A_i\,(1-\max_k E_{ik})$ & context-dep. & ref.\ + chunks\\
\code{faithfulness} & $\frac1p\sum_i \max_k E_{ik}$ & higher & chunks only\\
\bottomrule
\end{longtable}}

\code{precision}, \code{recall} and \code{f1} are the paper's \emph{overall} metrics. \code{claim_recall} and \code{context_precision} are \emph{retriever} metrics; they do not depend on the answer and are therefore identical for R1--R5. The remaining six are \emph{generator} diagnostics. Only \code{faithfulness} is defined without a reference: \code{NO_REFERENCE} mode computes it alone (\src{ragchecker.py:105}). Per-item values lie in $[0,1]$. The package's aggregate report multiplies the mean by 100 and rounds it (\src{evaluator.py:244}), but our tables report the per-item values.

\paragraph{LLM versus deterministic parts.}
\emph{LLM}: triplet extraction (response and reference) and the four entailment checks.
\emph{Deterministic}: triplet parsing (regex), joint-label parsing (regex for the three words), padding with Neutral, per-passage merging, and all eleven formulas above.

\subsection{C. What our code actually does}

\code{RAGCheckerAdapter.evaluate_many} (\src{ragchecker.py:86}):
\begin{enumerate}
\item \textbf{Inputs.} The adapter builds the official \code{RAGResults} payload: \code{query_id}=\code{item_id}, \code{query}, \code{gt_answer} = the golden reference text, \code{response} = the answer, and \code{retrieved_context} = \code{[{doc_id:"ctx0", text:ctx0}, {doc_id:"ctx1", text:ctx1}]}. The five answers to one question are sent together.
\item \textbf{Evaluator.} The adapter constructs the official class (\src{ragchecker.py:64}); apart from the hook, every argument is a package default:
\begin{lstlisting}
RAGChecker(extractor_name="custom/judge", checker_name="custom/judge",
           custom_llm_api_func=self._llm, joint_check=True,
           batch_size_extractor=32, batch_size_checker=32)
\end{lstlisting} \code{ev.evaluate(rr, all_metrics)} is then called unchanged.
\item \textbf{Triplet extraction (LLM).} The prompt \code{LLM_TRIPLET_EXTRACTION_PROMPT_Q} (\src{refchecker/extractor/extractor_prompts.py:178}) shows two in-context examples (Tesla Optimus, Andre Weiss) and ends with:
\begin{lstlisting}
Now generate the KG for the following candidate answer based on the provided question:
### Question:
What exact date was Alex Meron born?
### Candidate Answer:
Archive A gives 3 March 1990, while Registry B gives 17 March 1990. Both agree ...
### KG:
\end{lstlisting}
The judge returned the four triplets shown in Table~\ref{tab:rc-r5}. The reference is extracted with the same prompt. Because the judge cache is single-flight (Section~\ref{sec:cache}), the reference is extracted once and the same reference triplets are shared by all five answers.
\item \textbf{Joint checking (LLM).} \code{JOINT_CHECKING_PROMPT_Q} (\src{refchecker/checker/checker_prompts.py}) lists up to five triplets and asks for one label per line:
\begin{lstlisting}
If the claim is supported by ANY passage in the reference, answer 'Entailment'.
If NO passage in the reference entail the claim, and the claim is contradicted with some
passage in the reference, answer 'Contradiction'.
If NO passage entail or contradict with claim, or DOES NOT contain information to verify
the claim, answer 'Neutral'.
### Question: [QUESTION]
### Reference: [REFERENCE]
### Claims:
("Archive A", "gives Alex Meron's birth date as", "3 March 1990")
...
\end{lstlisting}
For \code{retrieved2*} checks, \emph{each chunk is a separate prompt}, so the judge sees one document at a time. For \code{answer2response}/\code{response2answer} the reference is the full text of the other side. Labels are parsed with a word-boundary regex for the three label words. If fewer labels than claims are found, RefChecker pads with Neutral silently (\src{llm_checker.py:144}). The adapter wraps the parser only to \emph{observe} this and record a warning (\src{ragchecker.py:69-80}); it never alters the labels.
\item \textbf{Scoring.} The \code{computation.py} functions run per item, and the adapter copies \code{r.metrics} into eleven \code{MetricValue}s. All four check matrices, both claim lists, the reference text and the reference id go into \code{raw}.
\end{enumerate}

\subsection{D. Worked examples (C1, judge \code{gpt-6-luna})}

\begin{table}[h]
\centering\small
\begin{tabularx}{\textwidth}{lXcc}
\toprule
$i$ & Response triplet $m_i$ (R5, Q\_DATE) & $(E_{i,A},E_{i,B})$ & $A_i$\\
\midrule
1 & (Archive A, gives Alex Meron's birth date as, 3 March 1990) & (Ent, Neu) & Ent\\
2 & (Registry B, gives Alex Meron's birth date as, 17 March 1990) & (Neu, Ent) & Ent\\
3 & (Archive A and Registry B, agree on \ldots\ birth month and year, March 1990) & (Ent, Neu) & Ent\\
4 & (Alex Meron's precise birth day, is, disputed) & (Neu, Neu) & Ent\\
\midrule
$j$ & Reference triplet $g_j$ (shared by R1--R5) & $(G_{j,A},G_{j,B})$ & $B_j$ (R5)\\
\midrule
1 & (Archive A, reports \ldots\ birth date as, 3 March 1990) & (Ent, Neu) & Ent\\
2 & (Registry B, reports \ldots\ birth date as, 17 March 1990) & (Neu, Ent) & Ent\\
3 & (Archive A and Registry B, agree on \ldots, March 1990) & (Neu, Neu) & Ent\\
4 & (Available evidence, does not establish \ldots\ correct birth day, unknown) & (Neu, Con) & Ent\\
\bottomrule
\end{tabularx}
\caption{RAGChecker raw judgments for \code{C1|Q_DATE|R5}, from \code{raw} in \src{results.jsonl}.}
\label{tab:rc-r5}
\end{table}

\paragraph{Why R1 has faithfulness 1.00 on Q\_DATE although Registry B contradicts it.}
R1 yields one triplet $m_1$ = (Alex Meron, date of birth, 3 March 1990), with $\code{retrieved2response}=[[\text{Ent},\text{Con}]]$. Then
\[
\max_k E_{1k} = \max(1,0)=1 \quad\Rightarrow\quad \code{faithfulness}=\tfrac11\cdot 1 = 1.00 .
\]
The formula takes the \emph{maximum} over chunks: one entailing chunk makes a claim faithful, and a contradiction from another chunk has no effect. In addition, \code{to_bool} would map the Contradiction to 0 even if it were counted. The checker prompt applies the same rule (``supported by ANY passage $\Rightarrow$ Entailment''). By construction, faithfulness cannot register cross-document conflict. The other values for R1 on Q\_DATE follow from the same matrices: $A=[\text{Ent}]$ gives precision $1.00$, and $B=[\text{Neu}]\times4$ gives recall 0 and therefore $\code{f1}=0$.

\paragraph{Why R5 has faithfulness 0.75 on Q\_DATE.}
From Table~\ref{tab:rc-r5}, $\max_k E_{ik}=(1,1,1,0)$, so
\[
\code{faithfulness}=\tfrac14(1+1+1+0)=0.75 .
\]
The claim ``the precise day is disputed'' was labelled Neutral against each document separately. Neither document alone states a dispute: the dispute exists only \emph{between} the two documents, and RAGChecker never shows the judge both chunks together. This is the structural reason the most conflict-aware answer scores below R1--R3. Triplet~3 (``A and B agree on March 1990'') shows that judge leniency matters as well. It was labelled Entailment from Archive A alone, although that document says nothing about Registry B. A stricter judge would have labelled it Neutral, giving $0.50$; this is the value Claude Sonnet produced in the cross-judge check.
Other R5 values: $A_i=1$ for all $i$ gives precision $1.00$, and $B_j=1$ for all $j$ gives recall $1.00$ and F1 $1.00$. Self-knowledge is $\frac14\sum_i A_i(1-\max_kE_{ik})=\frac14(0+0+0+1)=0.25$ (claim 4 is ``correct'' per the reference but in no chunk). Hallucination is 0.
The answer-independent retriever metrics are the same for every answer: $\max_k G_{jk}=(1,1,0,0)$ gives \code{claim_recall} $=0.50$, and $\max_j G_{jk}=(1,1)$ gives \code{context_precision} $=1.00$.

\paragraph{A judge-dependent precision value: R4 on Q\_DATE.}
R4's triplet (Alex Meron, date of birth, uncertain) received $A_1=$ Entailment from the reference, which says ``the available evidence does not establish which day is correct''. So precision is $1.00$ for the vaguest answer. $E_1=(\text{Con},\text{Con})$ gives faithfulness 0, and therefore $\code{self_knowledge}=1\cdot(1-0)=1.00$. Under both Claude judges R1's precision changed from $1.00$ to $0.00$ (\src{judge_comparison_agreement.csv}). The reference-based overall metrics clearly depend on how the judge reads a hedged reference.

\paragraph{Extractor-dependent claim counts: R1 on Q\_YEAR.}
The extractor produced two triplets from the single R1 sentence: (date of birth, 3 March 1990) and (year of birth, 1990). Against the reference ``born in 1990'', $A=(\text{Neu},\text{Ent})$, so precision is $1/2$. $m_1$ is incorrect but entailed by a relevant chunk, so \code{noise_sensitivity_in_relevant} $=1/2$. A different extraction of the same sentence would give different values.

%=====================================================================
\section{Q-CARE}
\label{sec:qc}

\subsection{A. Original paper and benchmark}

\begin{description}[leftmargin=0pt]
\item[Paper.] Jeonghwan Choi, Taewon Yun, Minjeong Ban, Gyeonghun Sun, Jae-Gil Lee, Hwanjun Song (KAIST). \emph{Towards Query-Agnostic RAG Evaluation via Query Coverage and Claim Verifiability.} Conference on Language Modeling (COLM) 2026. arXiv:\href{https://arxiv.org/abs/2608.11238}{2608.11238}.
\item[Code and data.] \url{https://github.com/DISL-Lab/Q-CaRE-COLM-26} (MIT), pinned at \code{e8b48e1}; dataset at \url{https://huggingface.co/datasets/DISLab/Q-CARE}.
\item[Contribution type.] There are three contributions: a \emph{reference-free evaluation framework} with three generator and two retriever metrics; a \emph{benchmark} (testbed); and a \emph{human-agreement study} comparing Q-CARE with existing RAG metrics, including RAGChecker.
\item[Benchmark.] 800 queries, 100 from each of 8 datasets. Close-ended: NQ, NewsQA, HotpotQA, FinQA. Open-ended: PubMedQA, LoTTE-Science, LoTTE-Technology, ELI5. Each record stores the top-30 BM25 and ANCE chunks and answers from eight RAG systems (GPT-5, Claude-Sonnet, Gemini-2.5-pro, GPT-oss-20B, Qwen3-30B, Qwen3-4B, Gemma3-27B, Gemma3-4B). The pipeline scores the top 10 chunks. Gold answers ship with the data but are never read.
\item[Judge models.] \emph{Repository default}: \code{Qwen/Qwen3-30B-A3B-Instruct-2507}, loaded locally with \code{transformers}, greedy decoding, \code{max_new_tokens=2000}, seed 1234 (\src{qcare/backbone.py}, \src{docs/REPRODUCE.md}). \emph{Paper}: per the README, the paper evaluates Qwen3-30B, Qwen3-80B, Llama3.1-8B and Llama3.3-70B as backbones. The README states that the backbone ``must be an open-weight chat model on HuggingFace; API-only models cannot be passed here''.
\end{description}

\paragraph{How our experiment differs.}
The original benchmark has 10 retrieved chunks per query from a real retriever, real multi-aspect questions and generated answers. We have two chunks, a single-focus question, fixed answers and a judge model reached through a CLI.

\subsection{B. Metric definitions and scoring logic}

Let $Q=\{q_1..q_s\}$ be the sub-queries and $F=\{f_1..f_t\}$ the atomic facts, both produced by the LLM. The LLM also provides four judgment sets:
$\mathrm{QFR}(q)\subseteq F$ (facts \emph{necessary} for $q$), $\mathrm{CFR}(k)\subseteq F$ (facts verified by chunk $k$), $\mathrm{QCC}(k)\subseteq Q$ (sub-queries fully answered by chunk $k$), and $\mathrm{QFC}(q)\in\{\text{Covered},\text{Not covered}\}$. The function \code{compute_metrics} (\src{qcare/metrics.py:84}) then computes
\begin{align}
\text{Completeness} &= \frac{|\{q: \mathrm{QFC}(q)=\text{Covered}\}|}{s},\\
\text{Conciseness} &= \frac{\bigl|\bigcup_{q:\,\mathrm{QFC}(q)=\text{Covered}} \mathrm{QFR}(q)\bigr|}{t},\\
\text{Verifiableness} &= \frac{\bigl|\bigcup_{k} \mathrm{CFR}(k)\bigr|}{t},\\
r_k &= \frac{|\mathrm{QCC}(k)|}{s},\quad k=1..10\ \ (r_k=0 \text{ for missing chunks}),\\
\text{C-Prec@10} &= \frac{1}{10}\sum_{k=1}^{10} r_k,\qquad
\text{C-nDCG@10} = \frac{\sum_{k} r_k/\log_2(k+1)}{\sum_k r^{\downarrow}_k/\log_2(k+1)} .
\end{align}
All five metrics are higher-is-better. The function \code{_ratio} returns 0 both when the denominator is 0 and when the ratio exceeds 1 (\src{metrics.py:71}); the second case can occur under \code{parsing="paper"}, where a fact named twice is counted twice. \code{QFC} is queried only when $\mathrm{QFR}(q)\neq\emptyset$; otherwise the sub-query is set to Not covered without an LLM call (\src{pipeline.py:139-148}).

\begin{table}[h]
\centering\small
\begin{tabular}{lllll}
\toprule
Metric & Kind & Uses evidence & Uses gold & Depends on answer\\
\midrule
completeness & answer quality & no & no & yes\\
conciseness & answer quality & no & no & yes\\
verifiableness & answer quality & yes & no & yes\\
\code{precision_at_10} & retrieval & yes & no & no\\
\code{ndcg_at_10} & retrieval & yes & no & no\\
\code{num_atomic_facts} & diagnostic & no & no & yes\\
\bottomrule
\end{tabular}
\caption{Q-CARE quantities stored by \src{evaluators/qcare.py}.}
\end{table}

\paragraph{LLM versus deterministic parts.}
\emph{LLM}: query decomposition, atomic-fact extraction, and the four judgment types (one call per (sub-query, fact), per chunk, per (sub-query, chunk) and per sub-query).
\emph{Deterministic}: JSON parsing (\src{parsing.py}); yes/no parsing with \code{is_yes} (case-insensitive substring \code{"yes"}); set unions; ratios; nDCG.

\subsection{C. What our code actually does}

\code{QCAREAdapter} (\src{evaluators/qcare.py:54}) builds the official pipeline:
\begin{lstlisting}
QCAREPipeline(JudgeBackbone(judge),
              PromptLibrary(str(REPO / "configs/prompts.yaml")), parsing="paper")
\end{lstlisting}
For each item it calls \code{pipeline.evaluate(item_id, query, answer, contexts)} (\src{qcare/pipeline.py:201}) in a thread pool. \code{JudgeBackbone.generate(messages)} forwards the official \emph{system + user} messages to \code{judge.chat}. The Codex judge then places the Q-CARE system message in Codex's \code{base_instructions} slot (\src{codex.py:160}).

\begin{enumerate}
\item \textbf{Query decomposition.} The \code{query_decomposition} prompt asks the judge to classify the query type and return single-focus queries ``EXACTLY as is''. For Q\_DATE it returned \code{{"Core subquery1": "What exact date was Alex Meron born?"}}, so $s=1$.
\item \textbf{Atomic facts.} The \code{atomic_fact_extraction} prompt asks for self-contained facts, keeps related information together, and requires at least one fact. For R5 the judge returned three facts (Table~\ref{tab:qc-r5}).
\item \textbf{Four judgments.} \code{query_fact_relevance}: ``Does the given ATOMIC FACT provide NECESSARY information to answer the SUB QUERY?'' (Yes/No, one call per fact). \code{chunk_fact_relevance}: one call per chunk listing all facts and returning \code{{"verified_facts": [...]}}. \code{query_chunk_coverage}: ``Does the given CHUNK provide COMPLETE information to FULLY answer the SUBQUERY?''. \code{query_fact_coverage}: whether the selected facts together cover the sub-query.
\item \textbf{Scoring and storage.} The official \code{compute_metrics}. The adapter stores the five metrics plus \code{num_atomic_facts}, and \code{raw} holds \code{decomposed_query}, \code{atomic_facts}, \code{relevance_check} and the official \code{metrics} block, including the counts.
\end{enumerate}

\subsection{D. Worked examples (C1, judge \code{gpt-6-luna})}

\begin{table}[h]
\centering\small
\begin{tabularx}{\textwidth}{lXccc}
\toprule
Key & Atomic fact (R5, Q\_DATE) & QFR($q_1$) & Chunk 1 & Chunk 2\\
\midrule
fact1 & Archive A gives Alex Meron's birth date as 3 March 1990. & yes & yes & --\\
fact2 & Registry B gives Alex Meron's birth date as 17 March 1990. & yes & -- & yes\\
fact3 & Archive A and Registry B agree that Alex Meron was born in March 1990, but they disagree on the precise day. & no & -- & --\\
\bottomrule
\end{tabularx}
\caption{Q-CARE judgments for \code{C1|Q_DATE|R5}. $\mathrm{QFC}(q_1)=$ Covered; $\mathrm{QCC}(1)=\mathrm{QCC}(2)=\{q_1\}$.}
\label{tab:qc-r5}
\end{table}

\paragraph{Why R5 has verifiableness 0.67.}
$\bigcup_k \mathrm{CFR}(k)=\{\text{fact1}\}\cup\{\text{fact2}\}=\{\text{fact1},\text{fact2}\}$ and $t=3$, so
\[
\text{Verifiableness}=\frac{2}{3}=0.67 .
\]
Fact~3 compares the two sources. \code{chunk_fact_relevance} is asked one chunk at a time, and no single chunk can verify a statement about what \emph{another} chunk says. This is the same structural effect as R5's claim~4 in RAGChecker. How many facts the extractor produces also matters. Under Claude Sonnet and Opus the extractor produced four facts for R5, two verified, giving $0.50$ (\src{judge_comparison_agreement.csv}). On Q\_YEAR R5 again has three facts, with fact3 unverified, giving $2/3$.

\paragraph{Why R5 has conciseness 0.67 on Q\_DATE but 1.00 on Q\_YEAR.}
On Q\_DATE the judge marked fact3 as not \emph{necessary} for ``What exact date\ldots?'', so the covered union is $\{\text{fact1},\text{fact2}\}$ and Conciseness $=2/3$. On Q\_YEAR it marked all three facts necessary, giving $3/3$. Completeness is $1/1=1.00$ in both cases.

\paragraph{Why R2 and R3 receive completeness 0 but verifiableness 1 on Q\_DATE.}
For R2, $\mathrm{QFR}(q_1)=\emptyset$: ``born in March 1990'' was judged not necessary for an exact-date question. The pipeline therefore records $\mathrm{QFC}(q_1)$ as Not covered without asking, so completeness and conciseness are 0/1 and 0/1. Both chunks verify the fact, so verifiableness is $1/1$. On Q\_DATE, R1 is ``complete'' while R5 is ``complete'' but less ``concise''.

\paragraph{Why every answer has C-Prec@10 $=0.20$ and C-nDCG@10 $=1.00$.}
Both chunks fully answer $q_1$, so $r=(1,1,0,\dots,0)$, giving $\text{C-Prec@10}=(1+1)/10=0.20$. This is the maximum possible with two chunks. The ranking is already ideal, so the DCG equals the IDCG ($10/\log_2 2 + 10/\log_2 3$ in both; the factor 10 is a no-op kept by upstream) and nDCG is $1.00$. The retriever metrics give no information about the answers and do not register the conflict: each chunk ``completely answers'' the exact-date question on its own.

%=====================================================================
\section{C1 results at a glance}
\label{sec:results}

\begin{table}[h]
\centering\small
\begin{tabular}{llccccc}
\toprule
Evaluator & Metric & R1 & R2 & R3 & R4 & R5\\
\midrule
ConflictScore & CS-C\,/\,CS-R ($\downarrow$) & 1\,/\,.50 & 0\,/\,0 & 0\,/\,0 & 0\,/\,1 & 0\,/\,0\\
\midrule
RC & precision & 1.00 & 1.00 & 1.00 & 1.00 & 1.00\\
RC & recall & 0.00 & 0.25 & 0.00 & 0.25 & 1.00\\
RC & f1 & 0.00 & 0.40 & 0.00 & 0.40 & 1.00\\
RC & claim\_recall\,/\,context\_precision & \multicolumn{5}{c}{0.50\,/\,1.00 for all answers}\\
RC & context\_utilization & 0.00 & 0.00 & 0.00 & 0.00 & 1.00\\
RC & self\_knowledge & 0.00 & 0.00 & 0.00 & 1.00 & 0.25\\
RC & faithfulness & 1.00 & 1.00 & 1.00 & 0.00 & 0.75\\
RC & hallucination, noise sens. (rel., irr.) & \multicolumn{5}{c}{0.00 for all answers}\\
\midrule
QC & completeness & 1.00 & 0.00 & 0.00 & 0.00 & 1.00\\
QC & conciseness & 1.00 & 0.00 & 0.00 & 0.00 & 0.67\\
QC & verifiableness & 1.00 & 1.00 & 1.00 & 0.00 & 0.67\\
QC & C-Prec@10\,/\,C-nDCG@10 & \multicolumn{5}{c}{0.20\,/\,1.00 for all answers}\\
\bottomrule
\end{tabular}
\caption{C1, Q\_DATE, judge \code{codex:gpt-6-luna:low} (\src{outputs/runs/main/results.jsonl}).}
\end{table}

\begin{table}[h]
\centering\small
\begin{tabular}{llccccc}
\toprule
Evaluator & Metric & R1 & R2 & R3 & R4 & R5\\
\midrule
ConflictScore & CS-C\,/\,CS-R ($\downarrow$) & 1\,/\,.50 & 0\,/\,0 & 0\,/\,0 & 0\,/\,1 & 0\,/\,0\\
RC & precision\,/\,recall\,/\,f1 & .50/1/.67 & 1/1/1 & 1/1/1 & 0/0/0 & 0/1/0\\
RC & noise\_sensitivity\_in\_relevant & 0.50 & 0.00 & 0.00 & 0.00 & 1.00\\
RC & hallucination & 0.00 & 0.00 & 0.00 & 1.00 & 0.00\\
RC & faithfulness & 1.00 & 1.00 & 1.00 & 0.00 & 1.00\\
QC & completeness\,/\,conciseness & 1/1 & 1/1 & 1/1 & 0/0 & 1/1\\
QC & verifiableness & 1.00 & 1.00 & 1.00 & 0.00 & 0.67\\
\bottomrule
\end{tabular}
\caption{C1, Q\_YEAR, same run. RAGChecker R5 faithfulness is 1.00 here because the extractor produced four attribution triplets (e.g.\ (Archive A, states Alex Meron was born in, 1990)), each entailed by one chunk, and no ``disputed'' triplet.}
\end{table}

%=====================================================================
\section{Experimental adaptations}
\label{sec:adapt}

\subsection{Judge substitution and model identifiers}

\begin{table}[h]
\centering\small
\begin{tabularx}{\textwidth}{p{2.4cm}>{\raggedright\arraybackslash}X}
\toprule
Evaluator & Original judge (paper / repository)\\
\midrule
ConflictScore & Paper: Gemini-3.5-Flash (main results), GPT-5.4-mini and Claude-haiku-4.5 (App.~B), GPT-4.1 for claim transformation. Repository: \code{gpt-4.1}, temperature 0.\\
RAGChecker & Paper: Llama3-70B-Instruct (extractor and checker). Code default: \code{bedrock/meta.llama3-70b-instruct-v1:0}; README examples: \code{bedrock/meta.llama3-1-70b-instruct-v1:0}.\\
Q-CARE & Repository: \code{Qwen/Qwen3-30B-A3B-Instruct-2507}, local, greedy, seed 1234. Paper: four open backbones (per README).\\
\midrule
\multicolumn{2}{p{\dimexpr\textwidth-2\tabcolsep}}{\textbf{This project, all three evaluators.} Main judge: backend \code{codex-cli}, requested model \code{gpt-6-luna}, \code{model_reasoning_effort=low}, harness \code{codex-judge-v1}, replicate 0; judgments generated with Codex CLI \code{0.167.0-alpha.4.1} (CS-C/CS-R prompts added on 15 and 19 November: \code{0.169.0-alpha.11.1}; the 33 RAGChecker prompts of the revised C4/C5 \textsc{Date} references, 23 November: \code{0.171.0-alpha.13}). Cross-judge check (C1 only): backend \code{claude-cli}, models \code{claude-sonnet-5-5} and \code{claude-opus-5-5}, Claude Code 2.1.322 (CS-C/CS-R prompts: 2.1.329), CLI default effort. One judge serves every LLM role within an evaluator (extractor, checker, decomposer).}\\
\bottomrule
\end{tabularx}
\caption{Original versus substituted judge configurations.}\label{tab:judges}
\end{table}

\noindent The model identifier recorded for Codex calls is the \emph{requested} one and is labelled \textbf{unverified} in run metadata (\code{model_identity}). In every Codex judge record, \code{model_reported_by_cli} is null, so the CLI never confirmed which model actually served the request. For Claude calls the served model is read from \code{modelUsage} and was recorded as \code{claude-opus-5-5} on the uncached Opus calls.

\subsection{How the Codex harness differs from a plain API call}

\code{CodexJudge._cmd} (\src{codex.py:160}) runs \code{codex exec} with \code{--ephemeral}, \code{--ignore-user-config}, \code{--sandbox read-only}, an empty temporary working directory, web search disabled, 19 agent features disabled, and seven instruction blocks switched off. The evaluator's system message (Q-CARE), or otherwise a fixed \code{NEUTRAL_SYSTEM} text, replaces Codex's \code{base_instructions}. The evaluator prompt is passed on stdin. Two consequences cannot be removed: (i) temperature cannot be set; (ii) about 8.2--8.7k input tokens are sent per call, of which about 6.9k are cached harness context (see \code{usage} in \src{judge_records.jsonl}). The evaluator's prompt is therefore never the model's whole context, unlike in the published configurations. Upstream \code{max_completion_tokens}/\code{max_new_tokens} limits are not forwarded.

\subsection{Reference-answer handling}

Only RAGChecker uses a reference. \code{cli.cmd_run} selects \code{GOLDEN} mode whenever the instance has a reference (\src{cli.py:102}). The reference is the per-(scenario, question) text in \code{GOLDEN_REFERENCES}, fixed before scoring and identical for R1--R5. Its construction principles P1--P7 are documented in the comment block at \src{data_factory.py:66-75}. Each result records the reference id and text (for example \code{GOLD:C1:Q_DATE}) and labels it ``evidence-grounded reference, NOT verified ground truth''. Earlier modes (\code{REF_A}, \code{REF_B}, \code{SHARED_MONTH}, \code{CONFLICT_AWARE}, \code{NO_REFERENCE}) are kept for reproducibility but are not part of the final C1 tables. ConflictScore and Q-CARE never see a reference.

\subsection{Prompt and parsing preservation}

RAGChecker and Q-CARE load their prompt templates, parsers and scoring functions from the pinned upstream code at run time: RAGChecker uses the installed package via its documented hook, Q-CARE uses \code{QCAREPipeline} with the repository's \code{prompts.yaml} and \code{parsing="paper"}. ConflictScore uses the paper's Appendix~E prompts verbatim and a port of the official label parser. The adapters observe failure modes that upstream code hides (RefChecker's Neutral padding, Q-CARE exceptions, judge errors) and turn them into \lbl{ERROR} status or warnings rather than scores.

\subsection{Caching, deduplication and reproducibility}
\label{sec:cache}

\begin{itemize}
\item \textbf{Content-addressed cache.} The key is $\mathrm{SHA256}$ over \{harness, model, effort, system text, prompt, output schema, replicate\} (\code{_key}, \src{codex.py:147}). Each response is written to \src{cache/judge/<key>.json} with its prompt, usage, timestamp and CLI version.
\item \textbf{Single-flight deduplication.} A lock per key (\src{codex.py:181}) ensures that identical prompts issued concurrently share \emph{one} call, as a temperature-0 backend would. This is why the RAGChecker reference is extracted once for R1--R5.
\item \textbf{Replay.} \code{offline=True} turns a cache miss into an error, and \code{max_new_calls} caps uncached calls. The reported results were produced by \src{scripts/run_all.sh} with \code{--cache-only}: every judgment was read from the cache (generated by the judge calls of 11, 15, 19 and 23 November 2026) and no new call was made, as \src{judgment_provenance.csv} records. With the cache present, rerunning reproduces every number exactly. Without it, results may differ, because the judge is not deterministic.
\item \textbf{Replicates.} \code{--replicate n} adds $n$ to the key and so forces fresh samples. No replicate other than 0 appears in the results.
\end{itemize}

\subsection{Technical limitations relevant to the conclusions}

\begin{enumerate}
\item \textbf{One sample per judgment.} Without temperature control, each label is a single draw. Values such as $0.67$ versus $0.50$ (Q-CARE R5) or R1 precision $1.00$ versus $0.00$ (RAGChecker) changed between judges. Conclusions that rest on a single label difference are fragile.
\item \textbf{Claim extraction is part of the score.} In all three evaluators the denominator ($n$, $p$, $t$) and the claim wording come from the judge. Examples: Opus splitting R2 into two ConflictScore claims (same scores); RAGChecker producing two triplets from R1; Q-CARE producing three or four facts for R5.
\item \textbf{Document-at-a-time checking.} All three evaluators check claims against one chunk at a time and combine the results with max or union rules. A statement about the \emph{relationship} between sources (``disputed'', ``both agree'') can be verified only through judge leniency.
\item \textbf{Unsupported configurations.} Q-CARE's README excludes API-only backbones. RAGChecker assumes a single correct reference. ConflictScore's paper scores long generated reports, not short fixed answers. Our results are therefore \emph{adapted reproductions with a substituted judge}, not reproductions of the published numbers.
\item \textbf{Small, degenerate evidence.} With two chunks, Q-CARE's C-Prec@10 cannot exceed 0.2, and the retriever metrics of both RAGChecker and Q-CARE are constant across answers.
\end{enumerate}

%=====================================================================
\section{Deviations from the published methodology}
\label{sec:disc}

Each item was checked against the paper, the pinned official code (identical to upstream HEAD on 15 November 2026) and the stored judge records. The classification says whether the deviation comes from the upstream code or documentation, or is a deliberate adaptation of this project.

\begin{enumerate}[label=\textbf{D\arabic*}, ref=D\arabic*, leftmargin=2.2em]
\item\label{d:cs-names} \textbf{ConflictScore's CS-C/CS-R have no official code.} There is no CS-C/CS-R code upstream. We implement them from the paper (\code{published_cs}); tests reproduce the paper's Fig.~2 and Fig.~5.
\item\label{d:cs-protocol} \textbf{The repository implements other protocols.} Its \src{src/} pipeline is the ConflictBench conflict-detection protocol (one claim rewritten from a dataset gold answer), and its example script uses other prompts and a support ratio. Neither computes CS-C/CS-R; neither is used.
\item\label{d:cs-prompts} \textbf{Repository prompts differ from the paper.} Paper Table~11 has no question and no ``both $\Rightarrow$ CONTRADICTS'' rule; Table~10 matches no repository prompt. We use only the paper's prompts.
\item\label{d:cs-conflict} \textbf{Conflict rule.} The repository labels a claim conflicted when $S/(S+C)<1$ (\src{eval_all_calibration.py:156-164}); the paper (Sec.~3) requires $|D^+|>0$ and $|D^-|>0$. We use the paper's rule.
\item\label{d:cs-irr} \textbf{CS-R for claims without relevant evidence.} The paper averages $|D^-|/(|D^+|+|D^-|)$ over claims but does not define the $0/0$ case. We exclude such claims from CS-R and count them as non-conflicted in CS-C (as in the paper's Fig.~2). Counting them as 0 would change two reported values (C4 R5: $0.25\to0.125$; NC R5: $0.33\to0.25$) and no ranking.
\item\label{d:cs-parse} \textbf{Claim-list and label parsing are not published.} We accept both ways judges follow Table~10's ``starting with "Claims: "'' instruction and port the official label parser. In the reported runs no fallback was needed and every label was valid.
\item\label{d:cs-label} \textbf{Judge models.} The paper's main ConflictScore judge is Gemini-3.5-Flash; GPT-4.1 is the repository default and the paper's claim-transformation model.
\item\label{d:rc-model} \textbf{RAGChecker judge.} The paper, the code default and the README examples name different Llama~3 models; all three are listed in Table~\ref{tab:judges}.
\item\label{d:rc-ref} \textbf{No ground truth.} \emph{Deliberate adaptation.} Ten of eleven RAGChecker metrics need a reference; ours are evidence-grounded texts written before scoring. Every result carries a \code{reference_dependent} flag, and these metrics measure agreement with our reference, not factual accuracy.
\item\label{d:qc-backbone} \textbf{Q-CARE on an API backbone.} \emph{Deliberate adaptation.} The README excludes API-only models; the 2000-token cap and retry loop of \src{backbone.py} are bypassed. The cap was never binding (largest visible output 94 tokens) and no retried failure reached scoring. Prompts and parsing are unchanged.
\item\label{d:qc-yaml} \textbf{Q-CARE YAML mentions ``--parsing robust''.} Upstream documentation error; no effect, since we use \code{paper}.
\item\label{d:judge-id} \textbf{Codex served model not confirmed.} No Codex record reports the served model; the identifier is the \emph{requested} one and is labelled unverified.
\item\label{d:temp} \textbf{Temperature not enforceable.} \emph{Documented limitation.} Neither CLI exposes temperature; each label is one sample. Reproducibility comes from the content-addressed cache.
\end{enumerate}

\vspace{1em}
\noindent\textbf{Sources consulted.}
ConflictScore paper PDF (ACL Anthology 2026.emnlp-main.2589), RAGChecker arXiv 2408.08067v2 (HTML), Q-CARE arXiv listing 2608.11238 and repository README/REPRODUCE.md, and the pinned upstream code under \src{external/} and \src{site-packages/ragchecker}, \src{site-packages/refchecker}. Paper claims attributed to a README rather than the paper text are marked as such.

\end{document}
